\section*{ПРИЛОЖЕНИЯ}
\label{sec:appendices}
\addcontentsline{toc}{section}{ПРИЛОЖЕНИЯ}

\subsection*{Приложение А. Възпроизвеждане на резултатите}
\label{app:reproduce}

Изходният код е публично достъпен в \url{https://github.com/cpp-for-everything/ORM}.
Резултатите със SQLite и компилационните тестове в глава~\ref{sec:evaluation} са
получени с версия \code{a17f4e7}. Резултатите със сървърните хранилища са от
изпълнение \code{41870491654} на работния процес \code{Docker Integration Tests}
върху комит \code{61255de}.
Наборът за оценка се изгражда и изпълнява със следните команди, от корена на
хранилището:

\begin{lstlisting}[language=bash,caption={Изграждане и изпълнение на набора за оценка},label={lst:app-reproduce}]
cmake -S . -B build-security -DORM_ENABLE_SECURITY=ON -DORM_ENABLE_SQLITE=ON
cmake --build build-security
ctest --test-dir build-security -L security --output-on-failure
\end{lstlisting}

Компилационните тестове изискват GCC или Clang. За изпълнение с Clang
променливите на средата \code{CC} и \code{CXX} се задават съответно на
\code{clang-18} и \code{clang++-18} преди конфигурирането. Резултатите се записват
в хранилището, заедно с изходния код на набора.

Програмите за сървърните хранилища се изграждат, когато съответният конектор е
включен, например с \code{-DORM\_ENABLE\_POSTGRESQL=ON}, и изискват работещ сървър,
достъпен през променливите \code{ORM\_*} на интеграционните тестове. Най-лесният
начин е образът \code{test-all} с услугите на Docker Compose, както го изпълнява
работният процес \code{Docker Integration Tests}. Той качва записаните резултати
като артефакт \code{security-evidence}.

\subsection*{Приложение Б. Състав на набора за оценка}

\begin{table}[H]
\centering
\caption{Части на набора за оценка}
\label{tab:app-files}
\begin{tabularx}{\textwidth}{L{5.4cm}Y}
\toprule
\textbf{Част} & \textbf{Съдържание} \\
\midrule
компилационни тестове & 18 програми и общата им постановка \\
проверка на компилационните тестове & проверка на изхода и на диагностиката на всеки компилационен тест \\
корпус & 66 полезни товара с колони за идентификатор, клас, диалект, техника, текст, очакван ефект и източник \\
източници на корпуса & 11 публични източника \\
тест със SQLite & диференциалният тест със SQLite \\
тестове със сървъри & пет програми: диференциалните тестове със сървърите на PostgreSQL, MySQL, MongoDB, Redis и Cassandra \\
общ драйвер & общият драйвер на сървърните тестове: контрол, оракул, резюме и правило за неуспех \\
именувани тестове & четири именувани теста на пътя през заместител \\
програми за възпроизвеждане & двете програми, възпроизвеждащи пропуска преди поправката \\
записани резултати & резюмета, записи по полезни товари и диагностики на компилаторите \\
\bottomrule
\end{tabularx}
\end{table}

Листинг~\ref{lst:app-corpus} показва вида на корпуса.

\begin{lstlisting}[language={},deletekeywords={orm},caption={Заглавен ред и първи запис на корпуса (записът е пренесен)},label={lst:app-corpus}]
"id","class","dialect","technique","payload","expected_effect","source_id"
"SQL-001","sql","generic","tautology","' OR '1'='1'--",
    "bypasses WHERE clause, returns all rows","patt_sqli"
\end{lstlisting}

\subsection*{Приложение В. Терминологичен речник}

Таблица~\ref{tab:app-glossary} съпоставя термините, използвани в текста, с
английските им съответствия и с идентификаторите в изходния код.

\begin{table}[H]
\centering
\caption{Терминологичен речник}
\label{tab:app-glossary}
\begin{tabularx}{\textwidth}{L{4.4cm}L{4.4cm}Y}
\toprule
\textbf{Термин} & \textbf{На английски} & \textbf{В кода} \\
\midrule
инжекция & injection & \\
полезен товар & payload & колона \code{payload} в корпуса \\
екраниране & escaping & \code{to\_sql\_literal} \\
свързване на параметри & parameter binding & \code{sqlite3\_bind\_*}, \code{PQexecParams} \\
типизиран заместител & typed placeholder & \code{Placeholder<T>}, \code{ph<T, \_N>} \\
литерален операнд & literal operand & \code{is\_literal\_operand\_v} \\
междинно представяне & intermediate representation & \code{select\_query} и сродните \\
свойство & property & \code{property<T, "name">} \\
конектор & connector & \code{connector\_trait<DB>} \\
проверка на границата & boundary check & \code{check\_query\_shape} \\
контролен случай & control case & случаи 00, 14 и 15 \\
оракул & oracle & функцията \code{classify} \\
затворен отказ & fail closed & \\
объркване на параметрите & parameter confusion & \\
\bottomrule
\end{tabularx}
\end{table}
